\section{Two bundle-map lemmas}
\label{sec:bundle-map-lemmas}

We first construct maps whose collective images expand subspaces.  We
then normalize such a family by a fiberwise minimization.  Uniqueness of
the minimizer will make the normalization continuous on the base, even
when the bundles are nontrivial.

\begin{lemma}[Avoiding small collective images]
\label{lem:incidence}
Let $E,F$ be smooth Hermitian complex vector bundles of ranks $d,r$ over
a compact smooth manifold $Y$ of real dimension $2h$, and let $q\geq1$.
Let $\mathcal A$ be a finite set of integer pairs $(s,t)$ with
$1\leq s\leq d$ and $0\leq t\leq r$ such that
\begin{equation}
  (qs-t)(r-t)-s(d-s)>h
  \qquad ((s,t)\in\mathcal A).
  \label{eq:incidence-dimension}
\end{equation}
There are smooth bundle maps $T_1,\ldots,T_q:E\to F$ such that, for every
$y\in Y$, every $(s,t)\in\mathcal A$, and every complex $s$-plane
$S\subset E_y$,
\[
  \dim_{\mathbb C}\sum_{i=1}^q T_i(y)S>t.
\]
Here and below the sum denotes linear span.  The same statement holds
over a point, with $h=0$.
\end{lemma}

\begin{proof}
Choose a finite cover trivializing $\operatorname{Hom}(E,F)$ and a smooth
partition of unity subordinate to it, with each support contained in its
trivializing open set.  Multiply every local frame section by the
corresponding partition function and extend by zero.  Their complex
linear span is a finite-dimensional space $V$ of smooth bundle maps
whose evaluation map
\[
  V\longrightarrow\operatorname{Hom}(E_y,F_y)
\]
is surjective for every $y$: at least one partition function is nonzero
there, and its multiplied frame still spans the fiber.  Thus the space
$\mathcal V=V^q$, of some complex dimension $M$, parametrizes tuples with
surjective evaluation onto $\operatorname{Hom}(E_y,F_y)^q$.

Fix $(s,t)\in\mathcal A$.  The Grassmann bundle of triples $(y,S,U)$,
where $S\subset E_y$ is an $s$-plane and $U\subset F_y$ is a $t$-plane,
has real dimension
\[
  2h+2s(d-s)+2t(r-t).
\]
Over this bundle, evaluation followed by restriction and quotient gives
a smooth, fiberwise surjective complex-linear map from the trivial
bundle with fiber $\mathcal V$ to the bundle with fiber
$\operatorname{Hom}(S,F_y/U)^q$.  Surjectivity follows because any linear
map $S\to F_y/U$ extends to a map $E_y\to F_y$, and evaluation from $V$
is surjective.  The target has complex rank $qs(r-t)$, so the kernels
form a smooth vector bundle whose total space $\mathcal I_{s,t}$ has real
dimension
\[
  2M+2h+2s(d-s)+2t(r-t)-2qs(r-t)<2M,
\]
where the inequality is precisely \eqref{eq:incidence-dimension}.

The image of the smooth projection
$\mathcal I_{s,t}\to\mathcal V$ consists exactly of tuples having some
$s$-plane whose collective image lies in a $t$-plane.  A smooth image of a manifold of smaller dimension is a Lebesgue
null set, as in the lower-dimensional case of Sard's theorem
\cite{Sard}.  Indeed, on a compact coordinate patch the projection is
Lipschitz.  Subdivision into cubes of side $\delta$ covers its image
by $O(\delta^{-a})$ balls of radius $O(\delta)$, where $a<2M$ is
the domain dimension; their total $2M$-dimensional volume tends to
zero.  A countable exhaustion by such patches handles the possibly
noncompact total space.  The finite
union of these null sets over $\mathcal A$ cannot exhaust $\mathcal V$.
A tuple outside that union has the required properties.
\end{proof}

The next normalization uses a determinant potential. Gurvits's
operator-scaling theory uses determinant capacity to relate rank
nondecrease to approximate doubly stochastic scaling \cite[Section~4]{Gurvits2003}.
Related affine log-determinant functions are studied by Sra and
Hosseini \cite[Corollary~2.11]{SraHosseini2015}. We prove the version
needed here, including the strict convexity and continuity that allow
normalization over nontrivial bundles.

\begin{lemma}[Continuous normalization]
\label{lem:scaling}
Let $E,F$ be Hermitian complex vector bundles of positive ranks over a
compact manifold $Y$.  Suppose that continuous bundle maps
$T_1,\ldots,T_n:E\to F$ satisfy
\begin{equation}
  \dim_{\mathbb C}\sum_{i=1}^n T_i(y)S
  \geq 2\dim_{\mathbb C}S
  \qquad (y\in Y,\ S\subset E_y).
  \label{eq:expansion}
\end{equation}
Then there are continuous bundle maps $X_1,\ldots,X_n:E\to F$ with
\begin{equation}
  \sum_{i=1}^nX_i^*X_i=I_E,
  \qquad
  \sum_{i=1}^nX_iX_i^*\leq\frac23 I_F.
  \label{eq:scaled-bounds}
\end{equation}
\end{lemma}

\begin{proof}
Write $d=\operatorname{rank}E$ and $r=\operatorname{rank}F$.  In each
fiber define
\[
  D_y(P)=\sum_{i=1}^nT_i(y)P T_i(y)^*,
  \qquad
  g_y(P)=\log\det(I_{F_y}+D_y(P))-\frac32\log\det P
\]
for positive definite $P\in\operatorname{End}(E_y)$.  Both determinants
are intrinsic.  We will prove that $g_y$ has a unique minimizer depending
continuously on $y$.  The expansion factor $2$ and the coefficient
$3/2$ leave a strict margin that prevents this minimizer from escaping
to infinity.

\paragraph{Uniform compact sublevels.}
Order all eigenvalues decreasingly.  For an orthogonal rank-$s$
projection $p$ on $E_y$,
\[
  \operatorname{ran}D_y(p)=\sum_iT_i(y)\operatorname{ran}p.
\]
Indeed, both sides have orthogonal complement
$\bigcap_i\ker(pT_i(y)^*)$.  By \eqref{eq:expansion}, the $2s$-th
eigenvalue of $D_y(p)$ is positive; in particular $r\geq2d$.
Eigenvalues are continuous, and the Grassmann bundle of such projections
is compact.  Taking a minimum there and then over $1\leq s\leq d$ gives
$\epsilon>0$ such that
\[
  \mu_{2s}(D_y(p))\geq\epsilon
\]
for every $y,s,p$ under consideration.

If $P$ has eigenvalues $\lambda_1\geq\cdots\geq\lambda_d>0$, choose any
projection $p_s$ onto $s$ top eigenvectors.  Then $P\geq\lambda_s p_s$,
so positivity of $D_y$ and eigenvalue monotonicity give
\[
  \mu_{2s}(D_y(P))\geq\epsilon\lambda_s.
\]
Pairing the target eigenvalues with indices $2s-1,2s$, and discarding any
remaining nonnegative terms, yields
\begin{equation}
  \log\det(I+D_y(P))
  \geq2\sum_{s=1}^d\log(1+\epsilon\lambda_s).
  \label{eq:scaling-determinant-bound}
\end{equation}
No continuous choice of the projections $p_s$ is needed.

Set $f(t)=2\log(1+\epsilon t)-\frac32\log t$.  This function is bounded
below and tends to $+\infty$ at both endpoints of $(0,\infty)$.  Let
$b=\min f$ and $C=\max_{y\in Y}g_y(I_{E_y})$.  Equation
\eqref{eq:scaling-determinant-bound} shows that, whenever $g_y(P)\leq C$,
\[
  f(\lambda_j)\leq C-(d-1)b\qquad(1\leq j\leq d).
\]
Thus all eigenvalues of every such $P$ lie in a common compact interval
$[a,A]\subset(0,\infty)$.  Outside the compact set
$aI_{E_y}\leq P\leq AI_{E_y}$, we have
$g_y(P)>C\geq g_y(I_{E_y})$.  Minimizing on this set therefore gives a
global minimizer in each fiber, lying in the interior of the positive
definite cone.

\paragraph{Strict convexity along positive definite geodesics.}
Fix $P_0>0$ and a nonzero Hermitian $H$ in one source fiber, and put
\[
  P(u)=P_0^{1/2}\exp(uH)P_0^{1/2}.
\]
Choose an orthonormal eigenbasis $e_j$ for $H$, with eigenvalues $h_j$,
and set $v_{ij}=T_iP_0^{1/2}e_j$.  Then
\[
  D(P(u))=\sum_{i,j}e^{u h_j}v_{ij}v_{ij}^*.
\]
The Gram determinant expansion of $\det(I+CC^*)$, applied to the columns
$e^{u h_j/2}v_{ij}$, expresses $\det(I+D(P(u)))$ as a finite sum
$\sum_\alpha c_\alpha e^{u\alpha}$ with $c_\alpha\geq0$.  The empty
subset contributes $1$, and singleton subsets contribute
$\|v_{ij}\|^2e^{u h_j}$.  For every $j$, the total singleton coefficient
at $h_j$ is positive: the vector $P_0^{1/2}e_j$ is nonzero and
\eqref{eq:expansion} makes the common kernel of the maps $T_i$ trivial.
Since some $h_j\ne0$, the expansion contains positive terms at two
distinct exponents, $0$ and $h_j$.

The second derivative of the logarithm of a positive exponential sum is
the variance of its exponents under weights proportional to
$c_\alpha e^{u\alpha}$.  It is therefore strictly positive here.  On the
other hand,
\[
  \log\det P(u)=\log\det P_0+u\operatorname{tr}H
\]
is affine.  Hence $g(P(u))$ is strictly convex.  Any two distinct positive
definite matrices $P_0,P_1$ are joined by such a nonconstant curve, by
taking $H=\log(P_0^{-1/2}P_1P_0^{-1/2})$.  The minimizer is unique.

\paragraph{Continuity and gluing.}
In local unitary frames, $(y,P)\mapsto g_y(P)$ is continuous and the
minimizers lie in the common compact matrix set $aI\leq P\leq AI$.
If $y_k\to y$, every subsequential limit of their minimizers minimizes
$g_y$: pass to the limit in the minimizing inequality against any fixed
positive definite comparison matrix.  Uniqueness forces every such
limit to be the minimizer at $y$, proving continuity.  On overlaps of
unitary frames, $P$ and $I_{F_y}+D_y(P)$ change by unitary conjugation,
preserving both determinants.  Uniqueness
therefore makes the local minimizers glue to a continuous positive
definite endomorphism $P$ of $E$.

\paragraph{The normalization.}
Put $Q=I_F+D(P)$.  Differentiation at the minimizer in an arbitrary
Hermitian source direction $K$ gives
\[
  0=\operatorname{tr}\left[
     \left(\sum_iT_i^*Q^{-1}T_i-\frac32 P^{-1}\right)K
   \right].
\]
The trace pairing on Hermitian matrices is nondegenerate, so
$\sum_iT_i^*Q^{-1}T_i=\frac32P^{-1}$.  Define
\[
  X_i=\sqrt{\frac23}\,Q^{-1/2}T_iP^{1/2}.
\]
Continuous functional calculus makes these continuous bundle maps, and
the order of the factors gives
\begin{align*}
  \sum_iX_i^*X_i
  &=\frac23 P^{1/2}\left(\sum_iT_i^*Q^{-1}T_i\right)P^{1/2}
    =I_E,\\
  \sum_iX_iX_i^*
  &=\frac23 Q^{-1/2}D(P)Q^{-1/2}
    =\frac23(I_F-Q^{-1})\leq\frac23 I_F.
\end{align*}
These are \eqref{eq:scaled-bounds}.
\end{proof}
